\documentclass[11pt]{article}
\usepackage[letterpaper,margin=1in]{geometry}
\usepackage[T1]{fontenc}
\usepackage{lmodern,microtype}
\usepackage{amsmath,amssymb,amsthm,mathtools}
\usepackage[colorlinks=true,linkcolor=blue,citecolor=blue,urlcolor=blue]{hyperref}
\hypersetup{pdftitle={Asymptotically optimal graph signings},pdfauthor={}}
\numberwithin{equation}{section}
\newtheorem{theorem}{Theorem}[section]
\newtheorem{lemma}[theorem]{Lemma}
\newtheorem{proposition}[theorem]{Proposition}
\newtheorem{corollary}[theorem]{Corollary}
\theoremstyle{remark}
\newtheorem{remark}[theorem]{Remark}
\newcommand{\norm}[1]{\left\lVert#1\right\rVert}
\newcommand{\E}{\mathbb E}
\newcommand{\Prob}{\mathbb P}
\newcommand{\ind}{\mathbf 1}
\newcommand{\tr}{\operatorname{tr}}
\allowdisplaybreaks[1]
\setlength{\emergencystretch}{2em}
\title{Asymptotically optimal graph signings}
\author{}
\date{September 29, 2026}
\begin{document}
\maketitle
\begin{abstract}
We give a proposed proof that every finite simple graph of maximum
degree at most $d$ admits an edge signing whose signed adjacency
matrix has norm $(2+o_d(1))\sqrt d$, uniformly in the number of
vertices. More quantitatively, the argument gives
$\bigl(2+O(\log\log d/\log d)\bigr)\sqrt d$.
The proof uses a nonnegative determinant weight supported on
signings satisfying a prescribed spectral bound. Uniform lower
bounds on its vertex-deletion ratios imply that each fixed diagonal
moment, under the associated weighted measure, approaches the
corresponding moment of a simple-path tree. Polynomial minorants
then give simultaneous lower bounds on the two insertion Schur
forms. A coupled estimate for the determinant and its diagonal
cofactors closes the vertex induction at every radius greater than
$2\sqrt d$. The proof is self-contained and existential.
\end{abstract}

\begin{center}
\small\textit{Research draft. The proposed proof has not undergone
independent human review.}
\end{center}

\section{Introduction}

An edge signing of a finite simple graph $G=(V,E)$ is a map
$\sigma:E\to\{-1,1\}$. Its signed adjacency matrix is the real
symmetric matrix
\[
 (A_\sigma)_{uv}=
 \begin{cases}
 \sigma(uv),&uv\in E,\\
 0,&uv\notin E.
 \end{cases}
\]
In particular, its diagonal is zero. We write $\norm{\cdot}$ for
the Euclidean operator norm. The signing determines a $2$-lift of
$G$, and the eigenvalues newly introduced by the lift are exactly
the eigenvalues of $A_\sigma$.

Bilu and Linial~\cite{BL} conjectured that every $d$-regular graph
has a signing with norm at most $2\sqrt{d-1}$. They proved the
weaker bound $O(\sqrt{d\log^3d})$. Marcus, Spielman, and
Srivastava~\cite{MSS} obtained the sharp one-sided bound
$\lambda_{\max}(A_\sigma)\le2\sqrt{d-1}$; spectral symmetry gives
the two-sided bound for bipartite graphs. For general graphs,
Ravichandran and Srivastava~\cite{RS} obtained
$2\sqrt{2(d-1)}$, and Lin and Zhou~\cite{LZ} improved the coefficient
to $(3+\sqrt5)/2$. Xu~\cite{Xu} recently disproved the exact
conjecture by constructing a cubic graph for which every signing
has norm greater than $2\sqrt2$.

The approximate question asks whether the constant $2$ remains
valid asymptotically as the degree tends to infinity. The result
below addresses that question, including the essential requirement
that the error term be independent of graph order.

\begin{theorem}\label{thm:main}
For every $\gamma>0$ there is an integer $d_0(\gamma)$ such that,
for every integer $d\ge d_0(\gamma)$ and every finite simple graph
$G$ of maximum degree at most $d$, there is an edge signing
$\sigma$ satisfying
\[
 \norm{A_\sigma}<(2+\gamma)\sqrt d.
\]
\end{theorem}

\begin{theorem}\label{thm:rate}
There is an absolute constant $C$ such that, for every sufficiently
large integer $d$, every finite simple graph of maximum degree at
most $d$ has an edge signing $\sigma$ satisfying
\[
 \norm{A_\sigma}
 \le \left(2+C\frac{\log\log d}{\log d}\right)\sqrt d.
\]
\end{theorem}

The proof works directly with graphs of bounded maximum degree,
so it accommodates all vertex deletions without a regular
completion. No connectedness assumption or bound on $|V(G)|$ is
needed.

\subsection{Proof strategy}

Fix a radius $R>2\sqrt d$. We assign a signing the weight
$\det(I-A_\sigma^2/R^2)$ when its norm is less than $R$, and weight
zero otherwise. Let $Z_H$ be the average weight on a graph $H$.
Proving $Z_H>0$ proves existence of a desired signing.
Determinant weights and their insertion/deletion identities also
appear in the randomized repair method of
Jadbabaie, Saberi, and Sra~\cite[Section~6]{JSS}. We use those
elementary identities inside an induction on graph order.

The induction controls both $Z_H/Z_{H-v}$ and a normalized diagonal
cofactor. A lower bound on all deletion ratios transfers
concentration from independent incident signs to the
determinant-weighted measure. The formal Schur recursion for
diagonal moments then agrees with the recursion on the simple-path
tree. Since moment order decreases in this comparison, its
constants do not grow with the number of vertices.

The two resolvents at the insertion vertex are bounded below by
finite geometric polynomials. Moment comparison therefore
controls their lower tails even when the actual resolvents are
large. Finally, if the two positive insertion residuals
$\delta_+,\delta_-$ are at most $s$, the elementary inequality
\[
 \delta_++\delta_-\le s+\frac{\delta_+\delta_-}{s}
\]
couples the cofactor to the insertion determinant. Retaining the
product on the right is what permits the induction to close at
every $R>2\sqrt d$.

\subsection{Notation and parameters}

All graphs below are finite, simple, and undirected. For
$v\in V(H)$, write $H-v$ for the induced graph obtained by deleting
$v$, and $N_H(v)$ for its neighbor set. All unweighted expectations
over signings refer to independent uniform signs. The norm and
determinant of an empty matrix are taken to be $0$ and $1$,
respectively. For a real number $t$, put $t_+=\max(t,0)$.
The notation $B\preceq C$ means that $C-B$ is positive semidefinite.
All logarithms are natural.

Until Section~\ref{sec:rate}, fix a real number $0<x<1/4$ and an
integer degree bound $d\ge1$, and set
\begin{equation}\label{eq:parameters-basic}
 Q=d/x,\qquad R=\sqrt Q,\qquad
 a=\frac{1+\sqrt{1-4x}}2,\qquad \rho=2\sqrt x<1.
\end{equation}
Thus $a>1/2$ and $x=a(1-a)$. These parameters remain fixed when
vertices are deleted: they are not recomputed from the maximum
degree of an induced subgraph.
It suffices to prove positivity of $Z_H$ for every graph of maximum
degree at most $d$, once $d$ is sufficiently large in terms of $x$.

\section{Determinant weights and vertex insertion}\label{sec:weights}

For a graph $H$ and signing $\sigma$, put $T_H=A_\sigma(H)/R$ and
\[
 W_H(\sigma)=
 \begin{cases}
 \det(I-T_H^2),&\norm{T_H}<1,\\
 0,&\norm{T_H}\ge1.
 \end{cases}
 \qquad Z_H=\E_\sigma W_H(\sigma).
\]
We call a signing good when $\norm{T_H}<1$, and put $Z_\varnothing=1$.
For a vertex $v$ define
\[
 \mathcal J_H(v)=
 \E_\sigma\left[
 W_H(\sigma)\left((I-T_H)^{-1}_{vv}+(I+T_H)^{-1}_{vv}\right)
 \mathbf1_{\{\norm{T_H}<1\}}\right].
\]
The integrand is defined to be zero on nongood signings; inverses
there are not used. When $Z_H>0$, write $\nu_H$ for the probability
measure with density $W_H/Z_H$ relative to uniform signs.
For clarity, $Z_H$ and $\mathcal J_H(v)$ are defined even if
$Z_H=0$; the probability measure $\nu_H$ is used only after
positivity has been established.

Delete a vertex $v$, set $K=H-v$, and fix a good signing of $K$.
Let $U=T_K$ and let $c$ be the normalized incident column at $v$,
with entries $\pm Q^{-1/2}$ on $N_H(v)$ and zero elsewhere. Define
\[
 q_\pm=c^T(I\pm U)^{-1}c,\quad
 \delta_\pm=1-q_\pm,\quad
 \alpha=(\delta_+)_+(\delta_-)_+,\quad S=q_++q_-.
\]
Order $v$ first, so that
\[
 T_H=\begin{pmatrix}0&c^T\\c&U\end{pmatrix}.
\]
Because $I\pm U$ are positive definite, the extension is good
exactly when both $\delta_+$ and $\delta_-$ are positive. Also
$\det(I-T_H^2)=\det(I-T_H)\det(I+T_H)$.
The determinant and inverse Schur-complement formulas therefore give
\begin{align}
 W_H&=W_K\alpha,\nonumber\\
 W_H\left((I-T_H)^{-1}_{vv}+(I+T_H)^{-1}_{vv}\right)
 &=W_K(\delta_++\delta_-)
   \mathbf1_{\{\delta_+>0,\delta_->0\}}.
 \label{eq:star}
\end{align}
Indeed, on a good extension the two inverse diagonal entries
are $1/\delta_-$ and $1/\delta_+$, whose sum multiplied by
$\delta_+\delta_-$ is $\delta_++\delta_-$.
On a nongood extension both left-hand integrands are defined to be
zero. This includes the boundary $\norm{T_H}=1$; the strict
indicator in the second identity is essential.
If $K$ is not good, no extension is good, by interlacing.
Since $q_\pm\ge0$, we have $0\le\alpha\le1$, and therefore
\begin{equation}\label{eq:domination}
 W_H\le W_{H-v}.
\end{equation}
In particular, if $Z_H/Z_{H-v}\ge c_0>0$, then for every nonnegative
function $F$ of the core signing and incident signs,
\begin{equation}\label{eq:density}
 \E_{\nu_H}F
 \le c_0^{-1}\E_{\nu_{H-v}}\E_{\mathrm{star}}F.
\end{equation}
More precisely, relative to $\nu_{H-v}$ times uniform incident
signs, the density of $\nu_H$ is $\alpha/(Z_H/Z_{H-v})$.
We will use both the full domination in~\eqref{eq:density} and its
consequence for the marginal distribution of the core.

\section{The simple-path tree}\label{sec:path-tree}

Let $\mathcal T(H,v)$ be the tree whose vertices are the simple
paths $(v=v_0,v_1,\ldots,v_k)$ in $H$ starting at $v$. Two paths
are adjacent if one is obtained from the other by appending one
vertex at the end. The root is the path $(v)$, which is also denoted by
$v$ when indexing matrices. Removing the root leaves one component
for each $u\in N_H(v)$, rooted at $(v,u)$ and isomorphic to
$\mathcal T(H-v,u)$.
Its maximum degree is at most $d$, so its adjacency norm is at most
$2\sqrt d$. For completeness, orient a tree away from its root and
sum
$2|z_pz_w|\le z_p^2/\sqrt d+\sqrt d\,z_w^2$
over parent--child edges; every resulting coefficient is at most
$2\sqrt d$.

Write
\[
 a_j(H,v)=
 \left[\left(A_{\mathcal T(H,v)}/R\right)^j\right]_{vv},
 \quad
 g_H(v)=\sum_{j\ge0}a_j(H,v),
 \quad r_H(v)=1/g_H(v).
\]
The tree is bipartite, so $a_j=0$ for odd $j$, and
$0\le a_{2j}\le\rho^{2j}$. The series converges absolutely.
Root Schur complementation gives
\begin{equation}\label{eq:tree}
 g_H(v)=
 \frac1{1-Q^{-1}\sum_{u\in N_H(v)}g_{H-v}(u)}.
\end{equation}
For an isolated root, $g_H(v)=1$. If
$g_{H-v}(u)\le1/a$ for all neighbors, the denominator
in~\eqref{eq:tree} lies between $1-x/a=a$ and $1$.
Induction on the number of vertices consequently gives
\begin{equation}\label{eq:r-range}
 1\le g_H(v)\le1/a,\qquad a\le r_H(v)\le1.
\end{equation}
For $K=H-v$ set
\[
 h=\deg_H(v)/Q,\quad r=r_H(v),\quad q=1-r.
\]
Then
\begin{equation}\label{eq:q-reference}
 0\le h\le x,\qquad
 q=\frac1Q\sum_{u\in N_H(v)}g_K(u).
\end{equation}

\section{Moment comparison and lower resolvent estimates}\label{sec:moments}

\begin{lemma}\label{lem:moments}
Fix $0<x<1/4$ and $0<c_0\le1$. Suppose that a family of graphs is closed
under vertex deletion, has maximum degree at most $d$, and satisfies
$Z_F>0$ and $Z_F/Z_{F-u}\ge c_0$ for every graph and vertex in the
family. For every integer $j\ge0$ there is a constant $C_j=C_j(x,c_0)$,
independent of $d$, graph order, and the family, such that
\[
 \left\|(T_F^j)_{uu}-a_j(F,u)\right\|_{L^2(\nu_F)}
 \le \frac{C_j}{\sqrt Q}.
\]
\end{lemma}
\begin{proof}
We induct on $j$, proving the assertion simultaneously for all
graphs in the family and all choices of root. The cases $j=0,1$
are exact because both matrices have zero diagonal.
For a fixed graph $F$ and root $u$, put $K=F-u$ and
$h=\deg_F(u)/Q$. Let $c$ be the normalized incident column
at $u$, as in Section~\ref{sec:weights}, and define
\[
 b_\ell=c^TT_K^\ell c,\qquad
 \beta_\ell=\frac1Q\sum_{w\in N_F(u)}a_\ell(K,w).
\]
On good core signings, $|b_\ell|\le h\le x$ and
$|\beta_\ell|\le x$. Conditional on the core,
\[
 \E_{\mathrm{star}}b_\ell
 =\frac1Q\sum_{w\in N_F(u)}(T_K^\ell)_{ww},\qquad
 \operatorname{Var}_{\mathrm{star}}(b_\ell)
 \le\frac{2\deg_F(u)}{Q^2}\le\frac{2x}{Q}.
\]
Indeed, for a symmetric $k$ by $k$ matrix $M$ and independent signs
$s_i$, $\operatorname{Var}(s^TMs)=2\sum_{i\ne j}M_{ij}^2$.
The relevant principal submatrix of $T_K^\ell$ has norm at most one.
Write
$\bar b_\ell=Q^{-1}\sum_{w\in N_F(u)}(T_K^\ell)_{ww}$.
Full measure domination in~\eqref{eq:density} gives
\[
 \|b_\ell-\bar b_\ell\|_{L^2(\nu_F)}
 \le \sqrt{\frac{2x}{c_0Q}}.
\]
For $\ell\le j-2$, marginal domination, Minkowski's inequality,
and the lower-order hypothesis on $K$ give
\begin{align*}
 \|\bar b_\ell-\beta_\ell\|_{L^2(\nu_F)}
 &\le c_0^{-1/2}
 \left\|\frac1Q\sum_{w\in N_F(u)}
       \bigl((T_K^\ell)_{ww}-a_\ell(K,w)\bigr)
 \right\|_{L^2(\nu_K)}\\
 &\le\frac{\deg_F(u)}{Q\sqrt{c_0}}\,
       \frac{C_\ell}{\sqrt Q}
 \le\frac{xC_\ell}{\sqrt{c_0Q}}.
\end{align*}
Combining the last two estimates yields
\begin{equation}\label{eq:b-moment}
 \|b_\ell-\beta_\ell\|_{L^2(\nu_F)}
 \le\frac{\sqrt{2x}+xC_\ell}{\sqrt{c_0Q}}.
\end{equation}

Let $M_j=(T_F^j)_{uu}$. Formal Schur complementation at the root gives
\[
 \sum_{j\ge0}M_jz^j
 =\frac1{1-\sum_{\ell\ge0}b_\ell z^{\ell+2}},
 \qquad
 M_j=\sum_{\ell=0}^{j-2}b_\ell M_{j-2-\ell}\quad(j\ge2).
\]
Exactly the same recurrence holds for $a_j(F,u)$ with $b_\ell$
replaced by $\beta_\ell$, by the definition of the path tree.
Subtract the recurrences in the form
\[
 M_j-a_j
 =\sum_{\ell=0}^{j-2}
 \left[(b_\ell-\beta_\ell)M_{j-2-\ell}
       +\beta_\ell(M_{j-2-\ell}-a_{j-2-\ell})\right].
\]
On the support of $\nu_F$, $|M_t|\le1$. Thus the desired estimate
follows from~\eqref{eq:b-moment}, using only moment orders at most
$j-2$. For example, the finite recursive constants
\[
 C_0=C_1=0,\qquad
 C_j=(j-1)\sqrt{2x/c_0}
      +x(1+c_0^{-1/2})\sum_{\ell=0}^{j-2}C_\ell
\]
suffice. This is induction on moment order, not on the size of the
graph, which is why the constants are uniform in graph order.
\end{proof}

\begin{lemma}\label{lem:lower-q}
Fix $0<x<1/4$, $0<c_0\le1$, and $\eta>0$. For sufficiently large $d$,
the following holds for every graph $H$ of maximum degree at most
$d$ and every $v\in V(H)$. Put $K=H-v$ and suppose its induced
subgraphs form a family satisfying the hypotheses of
Lemma~\ref{lem:moments}. Sample the core from $\nu_K$ and the
incident signs independently, and define
$q_\pm=c^T(I\pm T_K)^{-1}c$ as in Section~\ref{sec:weights}.
Then
\[
 \Prob\{q_+<q-\eta\ \text{or}\ q_-<q-\eta\}\le\eta,
 \qquad q=1-r_H(v).
\]
\end{lemma}
\begin{proof}
Put $U=T_K$. For an integer $m\ge1$ define
\[
 P_{\pm,m}(U)=\sum_{j=0}^{2m-1}(\mp U)^j.
\]
For $\norm U<1$, the scalar geometric-series identity gives
\[
 0\preceq P_{\pm,m}(U)\preceq(I\pm U)^{-1},
 \qquad \norm{P_{\pm,m}(U)}\le2m.
\]
Consequently $q_\pm\ge\ell_\pm:=c^TP_{\pm,m}(U)c$.
Define the deterministic truncated reference
\[
 q_m=\frac1Q\sum_{u\in N_H(v)}
            \sum_{j=0}^{m-1}a_{2j}(K,u).
\]
Conditional on $U$, the quadratic-form variance estimate gives
\[
 \|\ell_\pm-\E_{\mathrm{star}}\ell_\pm
       \|_{L^2(\nu_K\times\mathrm{star})}
 \le \frac{2m\sqrt{2x}}{\sqrt Q}.
\]
The conditional mean is
$Q^{-1}\sum_{u\in N_H(v)}(P_{\pm,m}(U))_{uu}$.
All odd reference moments vanish, and the coefficients of
$P_{\pm,m}$ have absolute value one. Thus
Lemma~\ref{lem:moments} and Minkowski's inequality give
\[
 \|\E_{\mathrm{star}}\ell_\pm-q_m\|_{L^2(\nu_K)}
 \le\frac{x}{\sqrt Q}\sum_{j=0}^{2m-1}C_j.
\]
In particular,
\[
 \|\ell_\pm-q_m\|_{L^2(\nu_K\times\mathrm{star})}
 \le\frac{B_m}{\sqrt Q},
 \qquad
 B_m=2m\sqrt{2x}+x\sum_{j=0}^{2m-1}C_j.
\]
Also
\[
 0\le q-q_m\le\frac{x\rho^{2m}}{1-\rho^2}.
\]
Choose $m$ so this last bound is at most $\eta/2$. A union bound and
Chebyshev's inequality then bound the probability in the statement
by $8B_m^2/(Q\eta^2)$, which is at most $\eta$ for sufficiently large
$d$. All parameters are independent of the order of $H$.
\end{proof}

\section{Coupled insertion and cofactor estimates}\label{sec:coupling}

\begin{lemma}\label{lem:scalar-coupling}
Let $q_+,q_-\ge0$, and define
$\delta_\pm=1-q_\pm$, $\alpha=(\delta_+)_+(\delta_-)_+$, and
$D=(\delta_++\delta_-)\mathbf1_{\{\delta_+>0,\delta_->0\}}$.
Fix $q'\in[0,1)$ and put $s=1-q'$. If $q_\pm\ge q'$, then
\[
 1-\alpha\le s(q_++q_-)+(q')^2,\qquad
 D\le s+\alpha/s.
\]
If the lower bounds $q_\pm\ge q'$ hold only on an event $E$, then
pointwise
\[
 1-\alpha\le s(q_++q_-)+(q')^2+\mathbf1_{E^c},
 \qquad D\le s+\alpha/s+2\mathbf1_{E^c}.
\]
\end{lemma}
\begin{proof}
Fix $0\le q'<1$ and put $s=1-q'$. If $q_\pm\ge q'$, then
\begin{equation}\label{eq:loss}
 1-\alpha\le s(q_++q_-)+(q')^2.
\end{equation}
If both $q_\pm<1$, expand the product and use
$(q_+-q')(q_--q')\ge0$. If either is at least one, then
$q_++q_-\ge1+q'$ and the right side is at least one.
For the cofactor, on a good extension one has
$0\le\delta_\pm\le s$. Therefore
$(s-\delta_+)(s-\delta_-)\ge0$, or
\begin{equation}\label{eq:coupling}
 (\delta_++\delta_-)\mathbf1_{\{\delta_\pm>0\}}
 \le s+\alpha/s.
\end{equation}
This is also true on nongood extensions, where the left side is zero.
On any exceptional event where the assumption $q_\pm\ge q'$ is
unavailable, add one to the right side of~\eqref{eq:loss}, and two to
the right side of~\eqref{eq:coupling}. These additions suffice because
the loss is at most one and the good-extension cofactor is at most
two. In particular, there is no unbounded resolvent multiplied by an
exceptional-event indicator.
\end{proof}

\section{Proof of Theorem~\ref{thm:main}}\label{sec:induction}

\begin{proof}

Choose a number $\varepsilon>0$ such that
\[
 D_*:=a^2-a\varepsilon x>x.
\]
Such a choice exists because $a^2-x=a(2a-1)>0$.
For example $\varepsilon=(a^2-x)/(2ax)$ works.
Set $c_0=D_*/2$, and choose $\eta>0$ sufficiently small that, with
\[
 \kappa=\eta\varepsilon x+\eta^2+\eta,
\]
both
\begin{equation}\label{eq:parameters}
 \kappa\le D_*/2,\qquad
 \kappa(1/a+\varepsilon)+3\eta
 \le\varepsilon(D_*-x)
\end{equation}
hold. Finally choose $d$ large enough for Lemma~\ref{lem:lower-q}
with these fixed $x,c_0,\eta$.

We prove simultaneously for every graph $H$ of maximum degree at
most $d$ that
\begin{equation}\label{eq:invariants}
 Z_H>0,\qquad
 \frac{Z_H}{Z_{H-v}}\ge c_0,\qquad
 \frac{\mathcal J_H(v)}{Z_H}\le2g_H(v)+\varepsilon
 \quad(v\in V(H)).
\end{equation}
Use strong induction on the number of vertices, starting with the
empty graph. Fix $H,v$ and set $K=H-v$. Every induced subgraph of
$K$ satisfies the ratio hypothesis by induction, so
Lemma~\ref{lem:lower-q} applies. Let $\E_*$ denote expectation under
$\nu_K$ and independent incident signs. Set
\[
 q=1-r_H(v),\quad r=1-q,\quad
 q'=\max(q-\eta,0),\quad t=q-q'\in[0,\eta],\quad s=1-q'=r+t.
\]
The event $E=\{q_\pm\ge q'\}$ has probability at least $1-\eta$.
Moreover, by independent incident signs and the cofactor invariant
on $K$,
\begin{align}
 \E_*S
 &=\frac1Q\sum_{u\in N_H(v)}\frac{\mathcal J_K(u)}{Z_K}
 \nonumber\\
 &\le\frac1Q\sum_{u\in N_H(v)}(2g_K(u)+\varepsilon)
 =2q+\varepsilon h,\qquad h=\deg_H(v)/Q\le x.
 \label{eq:mean-S}
\end{align}
Using~\eqref{eq:loss} on $E$ and adding $\mathbf1_{E^c}$ elsewhere
gives
\begin{align*}
 f:=\frac{Z_H}{Z_K}=\E_*\alpha
 &\ge1-s(2q+\varepsilon h)-(q-t)^2-\eta\\
 &=r^2-r\varepsilon h-t\varepsilon h-t^2-\eta\\
 &\ge r^2-r\varepsilon h-\kappa.
\end{align*}
Since $D_*>x$ implies $\varepsilon x<a$, the function
$r\mapsto r^2-r\varepsilon x$ is increasing for $r\ge a$.
Thus
\begin{equation}\label{eq:f-lower}
 f\ge D_*-\kappa\ge c_0>0.
\end{equation}
This proves positivity and all required deletion-ratio bounds.

For the cofactor, apply~\eqref{eq:coupling} on $E$ and add
$2\mathbf1_{E^c}$ elsewhere. Using~\eqref{eq:star},
\[
 \frac{\mathcal J_H(v)}{Z_K}
 \le s+\frac f s+2\eta,
 \qquad
 \frac{\mathcal J_H(v)}{Z_H}
 \le\frac1r+\frac{r+3\eta}{f}.
\]
Consequently~\eqref{eq:f-lower} and~\eqref{eq:parameters} give
\begin{align*}
 \frac{\mathcal J_H(v)}{Z_H}-\frac2r
 &\le\frac{(r^2-f)/r+3\eta}{f}\\
 &\le\frac{\varepsilon h+\kappa/r+3\eta}{f}\\
 &\le\frac{\varepsilon x+\kappa/a+3\eta}{D_*-\kappa}
 \le\varepsilon.
\end{align*}
This closes~\eqref{eq:invariants}.
Since $v$ was arbitrary, the induction is complete.

Positive $Z_H$ implies that some signing has positive weight, and
therefore norm strictly less than $R=\sqrt{d/x}$.
Given $\gamma>0$, take $x=(2+\gamma)^{-2}<1/4$.
This proves Theorem~\ref{thm:main}.
\end{proof}


\section{A quantitative error term}\label{sec:rate}

\begin{proof}[Proof of Theorem~\ref{thm:rate}]
Put $s=\sqrt{1-4x}$ and restrict to $0<s\le1/2$. Thus
$x=(1-s^2)/4$, $a=(1+s)/2$, and $x=a(1-a)$.
Choose the auxiliary parameters exactly as follows:
\[
 \varepsilon=\frac{a^2-x}{2ax}=\frac{s}{2x},\qquad
 D_*=\frac a2,\qquad c_0=\frac a4\ge\frac18,\qquad
 \eta=\frac{s^2}{100}.
\]
These choices satisfy~\eqref{eq:parameters}. Indeed,
$\varepsilon x=s/2\le1/4$ and $\eta\le1/400$, so
\[
 \kappa=\eta(1+\varepsilon x)+\eta^2
 \le\frac{501}{400}\eta.
\]
Also $1/a\le2$, $\varepsilon\le4/3$, and consequently
\[
 \kappa(1/a+\varepsilon)+3\eta
 \le\frac{287}{40}\eta
 =\frac{287}{4000}s^2
 <\frac{s^2}{2}
 \le\frac{as^2}{4x}
 =\varepsilon(D_*-x).
\]
Finally, $\kappa\le501/160000<1/8\le D_*/2$.

We next make the degree threshold explicit. In
Lemma~\ref{lem:moments}, uniformly for $x\le1/4$ and $c_0\ge1/8$,
\[
 \sqrt{2x/c_0}\le2,\qquad
 x(1+c_0^{-1/2})\le\frac{1+\sqrt8}{4}<1.
\]
Hence the recursive constants satisfy
\[
 C_j\le2(j-1)+\sum_{\ell=0}^{j-2}C_\ell,\qquad C_0=C_1=0.
\]
Induction gives $C_j\le2^j$: using the bounds for lower indices, the
right side is at most $2(j-1)+2^{j-1}-1\le2^j$ for $j\ge2$.
Thus, for $m\ge1$,
\[
 B_m\le\sqrt2\,m+\frac{4^m-1}{4}\le4^m.
\]
For the last inequality, use $m\le4^{m-1}$ and $1+\sqrt2<4$.

Choose
\begin{equation}\label{eq:explicit-m}
 m=\left\lceil
       \frac{\log(100/s^4)}{-\log(1-s^2)}
    \right\rceil.
\end{equation}
Then $(4x)^m=(1-s^2)^m\le s^4/100$, so the reference tail in
Lemma~\ref{lem:lower-q} is at most
\[
 \frac{x(4x)^m}{1-4x}\le\frac{s^2}{400}<\frac{\eta}{2}.
\]
The exceptional probability is at most
$8\cdot16^m/(Q\eta^2)$, and is therefore at most $\eta$ whenever
\begin{equation}\label{eq:explicit-Q}
 Q\ge \frac{8\cdot10^6\,16^m}{s^6}.
\end{equation}
All the hypotheses of the preceding induction then hold.

Take $x=(2+\gamma)^{-2}$ with $\gamma>0$ sufficiently small.
Since $s^2\asymp\gamma$, equations~\eqref{eq:explicit-m}
and~\eqref{eq:explicit-Q} show that a sufficient degree threshold
can be chosen with
\begin{equation}\label{eq:d-threshold}
 \log d_0(\gamma)
 \le K\gamma^{-1}\log(1/\gamma)
\end{equation}
for an absolute constant $K$ and all sufficiently small $\gamma$.
Here one may take the integer ceiling of $x$ times the right side
of~\eqref{eq:explicit-Q} as the degree threshold.

Choose an absolute $C>K$ and set
$\gamma=C\log\log d/\log d$. For sufficiently large $d$,
$\log(1/\gamma)\le\log\log d$, and therefore
\[
 K\gamma^{-1}\log(1/\gamma)
 \le(K/C)\log d<\log d.
\]
Thus $d\ge d_0(\gamma)$, and Theorem~\ref{thm:main} gives the
claimed quantitative bound.
\end{proof}

\section{Scope and dependence of the estimates}

\subsection{Uniformity in graph order}
The moment constants recurse in moment degree, not graph order.
For fixed $x<1/4$, the path-tree series has a uniform geometric tail,
so only finitely many moments are needed. The single-vertex
partition-function ratios stay above a fixed $c_0>0$ throughout the
induction. The order of parameter choices is
$x$, then $\varepsilon,c_0$, then $\eta$, then the truncation degree,
then $d_0$. No parameter depends on the number of vertices.

\subsection{Role of the cofactor coupling}
Replacing~\eqref{eq:coupling} by the weaker bound
$\delta_++\delta_-\le2s$ would lose the term $\alpha/s$.
The resulting error recurrence has a factor of two too much near
the tree threshold. Retaining $\alpha$ permits the stability margin
$a^2-x>0$ for every $x<1/4$.



\subsection{Existence and limitations}

The argument proves positivity of a finite partition function; it
does not give an efficient method to sample the weighted measure
or compute its deletion ratios. Accordingly, the result here is
an existence theorem, with no polynomial-time algorithm asserted.
The degree threshold obtained from~\eqref{eq:explicit-Q} is very
large, and no effort has been made to optimize it.
The conclusion concerns the regime $d\to\infty$ and does not
assert the exact bound $2\sqrt{d-1}$ at fixed degree.

\begin{thebibliography}{9}

\bibitem{BL}
Y. Bilu and N. Linial,
\emph{Lifts, discrepancy and nearly optimal spectral gap},
Combinatorica \textbf{26} (2006), no.~5, 495--519.
\url{https://doi.org/10.1007/s00493-006-0029-7}.

\bibitem{JSS}
A. Jadbabaie, A. Saberi, and S. Sra,
\emph{Kadison--Singer partitions and Bilu--Linial graph signings
in polynomial time},
arXiv:2609.23855 (2026).
\url{https://arxiv.org/abs/2609.23855}.

\bibitem{LZ}
F. Lin and H. Zhou,
\emph{Improved bounds for the Bilu--Linial conjecture via spectral
recovery from mixed determinantal polynomials},
arXiv:2609.15715 (2026).
\url{https://arxiv.org/abs/2609.15715}.

\bibitem{MSS}
A. W. Marcus, D. A. Spielman, and N. Srivastava,
\emph{Interlacing families I: Bipartite Ramanujan graphs of all degrees},
Annals of Mathematics (2) \textbf{182} (2015), no.~1, 307--325.
\url{https://arxiv.org/abs/1304.4132}.

\bibitem{RS}
M. Ravichandran and N. Srivastava,
\emph{Asymptotically optimal multi-paving},
International Mathematics Research Notices \textbf{2021} (2021),
no.~14, 10908--10940.
\url{https://doi.org/10.1093/imrn/rnz111}.

\bibitem{Xu}
Z. Xu,
\emph{A $3$-regular counterexample to the Bilu--Linial signing
conjecture},
arXiv:2609.15591 (2026).
\url{https://arxiv.org/abs/2609.15591}.

\end{thebibliography}
\end{document}
